\documentclass[fleqn,compress,utf8,aspectratio=169,t]{beamer}
% use LMU theme
\usetheme{LMU}

% This setups a lot of recommended settings and packages you usually need
% anyway. So let me take care for you and your document stays clear.

% Examples: Tables, Graphics, Listings, Hyperref...
%
% Prevent slide breaks in the middle of a paragraph
\widowpenalties 1 10000
\raggedbottom

\makeatletter
\@ifundefined{KOMAClassName}{% if non-KOMA class
  \IfFileExists{parskip.sty}{%
    \RequirePackage{parskip}
  }{% else
    \setlength{\parindent}{0pt}
    \setlength{\parskip}{6pt plus 2pt minus 1pt}}
}{% if KOMA class
  \KOMAoptions{parskip=half}}
\makeatother

\IfFileExists{xurl.sty}{\usepackage{xurl}}{} % add URL line breaks if available
\IfFileExists{bookmark.sty}{\usepackage{bookmark}}{\usepackage{hyperref}}
\urlstyle{same} % disable monospaced font for URLs

\usepackage{listings}
\lstset{defaultdialect=[5.3]Lua}
\lstset{defaultdialect=[x86masm]Assembler}
\usepackage{longtable,booktabs}
\usepackage{caption}
% Make caption package work with longtable
\makeatletter
\def\fnum@table{\tablename~\thetable}
\makeatother
\usepackage{graphicx}
\makeatletter
\def\maxwidth{\ifdim\Gin@nat@width>\linewidth\linewidth\else\Gin@nat@width\fi}
\def\maxheight{\ifdim\Gin@nat@height>\textheight\textheight\else\Gin@nat@height\fi}
\makeatother
% Scale images if necessary, so that they will not overflow the page
% margins by default, and it is still possible to overwrite the defaults
% using explicit options in \includegraphics[width, height, ...]{}
\setkeys{Gin}{width=\maxwidth,height=\maxheight,keepaspectratio}
% Set default figure placement to htbp
\makeatletter
\def\fps@figure{htbp}
\makeatother
\usepackage[normalem]{ulem}
% Avoid problems with \sout in headers with hyperref
\pdfstringdefDisableCommands{\renewcommand{\sout}{}}
\setlength{\emergencystretch}{3em} % prevent overfull lines
%Tightlist Command
\providecommand{\tightlist}{%
  \setlength{\itemsep}{0pt}\setlength{\parskip}{0pt}}



\usepackage[style=ieee]{biblatex}
\addbibresource{lmu-example.bib}

\usepackage{lmodern}

%%%%%%%%%%%%%%%%%%%%%%%%%%%%%%%%%%%%%%%%%%%%%%%%%%%%%%%%%%%%%%%%%%%%%%%%%%%%%%%%
%%                                 Customizations                             %%
%%%%%%%%%%%%%%%%%%%%%%%%%%%%%%%%%%%%%%%%%%%%%%%%%%%%%%%%%%%%%%%%%%%%%%%%%%%%%%%%

% If you want no navigation uncomment this
%\beamertemplatenavigationsymbolsempty

% setup listings
\lstset{
  basicstyle=\ttfamily\color{black},
  showstringspaces=false
}

\lstdefinestyle{highlight}{
  keywordstyle=\color{red},
  commentstyle=\color{lmu@darkgray}
}

\lstdefinestyle{basetex}{
language={[LaTeX]Tex},
basicstyle=\color{black!40},
keywordstyle=\color{red!40},
commentstyle=\color{lmu@darkgray!40},
moredelim=**[il][\only<+>{\color{black}\lstset{style=highlight}}]{@}
}

\lstdefinestyle{basec}{
language=C,
basicstyle=\color{black!40},
keywordstyle=\color{red!40},
commentstyle=\color{green!40},
moredelim=**[il][\only<+>{\color{black}\lstset{style=highlight}}]{@}
}

%%%%%%%%%%%%%%%%%%%%%%%%%%%%%%%%%%%%%%%%%%%%%%%%%%%%%%%%%%%%%%%%%%%%%%%%%%%%%%%%
%%                                  Title Page Data                           %%
%%%%%%%%%%%%%%%%%%%%%%%%%%%%%%%%%%%%%%%%%%%%%%%%%%%%%%%%%%%%%%%%%
% helper command to add multiple authors
\newcommand{\newauthor}[2]{
\parbox[c]{0.26\textwidth}{
{\bfseries #1} \\
{\scriptsize{\href{mailto:#2}{#2}}}
}
%{#1}
}

% Authors
\author[Mustermann]{
  \newauthor{Max Mustermann\inst{1}}{mustermann@nm.ifi.lmu.de} \and
  \newauthor{Karl Wombat\inst{2}}{wombat@tum.de}
}

\institute[LMU]{
  \inst{1}%
  {MNM-Team, LMU München}
  \inst{2}%
  {TUM}
}

\date[\today]{\today}

\title{Title}
\subtitle{Subtitle}

\hypersetup{
  pdftitle={Title},
  pdfauthor={Author},
  hidelinks}


%%%%%%%%%%%%%%%%%%%%%%%%%%%%%%%%%%%%%%%%%%%%%%%%%%%%%%%%%%%%%%%%%%%%%%%%%%%%%%%%
%%                                  Document                                  %%
%%%%%%%%%%%%%%%%%%%%%%%%%%%%%%%%%%%%%%%%%%%%%%%%%%%%%%%%%%%%%%%%%%%%%%%%%%%%%%%%

\begin{document}

\begin{frame}
  \titlepage
\end{frame}

%%%%%%%%%%%%%%%%%%%%%%%%%%%%%%%%%%% Overview %%%%%%%%%%%%%%%%%%%%%%%%%%%%%%%%%%%

\section{Overview}

\subsection{Motivation}

\begin{frame}{Motivation}
  \begin{itemize}
    \item The LMU provides students, professors and other staff with
          \begin{itemize}
            \item Corporate Design Manual~\cite{lmu2023}
          \end{itemize}
          \pause
    \item The IT-Service for Physics provides some \LaTeX\ templates such as beamer
          theme~\cite{lmu2014}.

    \item The MNM Team again modernized this and changed the ratio to 16:9
  \end{itemize}
\end{frame}

\subsection{Usage}

\begin{frame}{Usage}
  \begin{itemize}
    \item install the style-files in your project-directory
          \begin{itemize}
            \item beamerthemeLMU.sty
            \item beamerinnerthemelmu.sty
            \item beamerouterthemelmu.sty
            \item beamercolorthemelmu.sty
            \item beamerfontthemelmu.sty
          \end{itemize}
  \end{itemize}
\end{frame}

\begin{frame}{Usage}
	\begin{itemize}
		\item install the logos project-directory
		\begin{itemize}
			\item lmulogogreen.png
			\item siegel.pdf
			\item mnmlogo.pdf
		\end{itemize}
		\item compile twice with pdflatex
	\end{itemize}
\end{frame}

\begin{frame}[fragile]{Usage}
  \begin{block}{Code}
    \begin{lstlisting}[style=basetex]
\documentclass[
    @compress,    % compresses navigation bar
]{beamer}

@\usetheme{LMU}   % enable LMU theme
\end{lstlisting}
  \end{block}
  \begin{itemize}
    \item<1-> navigation bar has to be compressed explicitly
    \item<2-> no options to \lstinline$\usetheme$ implemented yet
    \item<3-> any customization has to happen in the theme's code
  \end{itemize}
\end{frame}

\subsection{Documentation}

\begin{frame}{Documentation}
  At this time:
  \begin{itemize}
    \item<1-> no final documentation available
    \item<2-> documentation mostly as source code
    \item<3-> this very presentation serves as showcase
          %\item<4-> project page on our GitLab server \cite{Wan}
  \end{itemize}

\end{frame}

%%%%%%%%%%%%%%%%%%%%%%%%%%%%%%%%%%% Examples %%%%%%%%%%%%%%%%%%%%%%%%%%%%%%%%%%%

\section{Examples}

\subsection{Lists}

\begin{frame}{Lists}
  \begin{columns}[t]
    \column{.5\textwidth}
    \begin{itemize}
      \item This is a list
      \item using the
      \item itemize environment.
    \end{itemize}
    \column{.5\textwidth}
    \begin{enumerate}
      \item This is a list
      \item using the
      \item enumerate environment.
    \end{enumerate}
  \end{columns}
\end{frame}

\subsection{Formulas}

\begin{frame}{Formulas}
  Four formulas --- One Idea
  \begin{align}
    \int^{b}_{a} f'(x) \, dx                                 & = f(b) - f(a) \\
    \int^b_{a} \nabla f\,d\mathbf{l}                         & = f(b) - f(a) \\
    \int_{\mathcal{V}} \nabla\mathbf{f} \, dv                & =
    \oint_{\mathcal{S}} \mathbf{f} \, d\mathbf{a}                            \\
    \int_{\mathcal{S}} \nabla\times\mathbf{f} \, d\mathbf{a} & =
    \oint_{\mathcal{P}} \mathbf{f} \, d\mathbf{l}
  \end{align}
\end{frame}

\subsection{Blocks}

\begin{frame}{Blocks}
  \begin{columns}[t]
    \column{.5\textwidth}
    \begin{block}{Regular Block}
      Theorems, proofs and definitions will look the same.
    \end{block}
    \begin{theorem}
      Told you.
    \end{theorem}
    \begin{proof}
      See above.
    \end{proof}
    \column{.5\textwidth}
    \begin{definition}
      This is a definition.
    \end{definition}
    \begin{example}
      This is an example of an example.
    \end{example}
    \begin{alertblock}{Alert Block}
      Nothing to be \alert{alerted} about.
    \end{alertblock}
  \end{columns}
\end{frame}

\subsection{Figures}

\begin{frame}{Figures}
  Figures drawn with TikZ can be uncovered part by part:
  \begin{columns}[t]
    \column{.5\textwidth}
    \begin{itemize}
      \item<2-> Das
      \item<3-> ist
      \item<4-> das
      \item<5-> Haus
      \item<6-> vom
      \item<7-> Ni
      \item<8-> ko
      \item<9-> laus
    \end{itemize}
    \column{.5\textwidth}
    \begin{figure}
      \begin{tikzpicture}
        \useasboundingbox (-0.5,-.5) rectangle (3.5,4);
        \color{white}
        \fill (-0.5,-.5) rectangle (3.5,4);
        \color{red!75!black}
        \draw<2-> (0,0)     -- (0,2);
        \draw<3-> (0,2)     -- (1.5,3.5);
        \draw<4-> (1.5,3.5) -- (3,2);
        \draw<5-> (3,2)     -- (3,0);
        \draw<6-> (3,0)     -- (0,2);
        \draw<7-> (0,2)     -- (3,2);
        \draw<8-> (3,2)     -- (0,0);
        \draw<9-> (0,0)     -- (3,0);
      \end{tikzpicture}
    \end{figure}
  \end{columns}
\end{frame}

\subsection{Source Code}

\begin{frame}[fragile]{Source Code}
  Source Code may be highlighted part by part:
  \begin{block}{Code}
    \begin{lstlisting}[style=basec]
#include <sdtio.h>
@int main (void) {
@   printf("Hello World!\n");
@   return 0;
}
\end{lstlisting}
  \end{block}
\end{frame}

%%%%%%%%%%%%%%%%%%%%%%%%%%%%%%%%%%% Closing %%%%%%%%%%%%%%%%%%%%%%%%%%%%%%%%%%%%

\section{Closing}

\subsection{Outlook}

\begin{frame}{Outlook}
  The LMU Beamer Theme already provides you with a beautifully typeset template
  that follows the LMU corporate design. The following things are still missing:
  \begin{itemize}
    \item<2-> Incorporation of the LMU Fonts
    \item<3-> Customizability with Theme-Options
    \item<4-> More elaborate Color Theme
    \item<5-> Powerpoint-compatible Theme
    \item<6-> Portability to Article-Mode
    \item<7-> Packaging for Distributions
  \end{itemize}

\end{frame}

\subsection{Development}

\begin{frame}{Development}
  \begin{itemize}
    \item<1-> theme is still under development
    \item<2-> there are missing features
    \item<3-> there may be bugs
  \end{itemize}
  \uncover<4->{
    \begin{alertblock}{Caution}
      The LMU beamer theme is considered alpha!
    \end{alertblock}
  }
\end{frame}

\subsection{References}
\begin{frame}[allowframebreaks]
  \frametitle{References}
  %\nocite{*}
  \printbibliography[heading=none]
\end{frame}

\end{document}
